\chapter{Newton Solvers and Damping}

The default 3D PMG benchmark uses the same nonlinear idea as MATLAB:
\begin{itemize}
\item build residual and tangent information,
\item regularize the tangent,
\item solve for a correction,
\item line-search the update,
\item accept or reject the iteration.
\end{itemize}

What changed is that the current path is now explicitly tied to the reduced
free-space operator produced by owned-row distributed assembly.

\section{Current Mainline Newton Branch}

For the committed benchmark, the Newton decision tree collapses to one branch:
\begin{itemize}
\item the backend is \code{pmg\_shell},
\item \code{pmg\_shell} enforces
\code{full\_system\_preconditioner = false},
\item Newton therefore uses the reduced free operator path,
\item the owned-row tangent machinery stays underneath the builder,
\item the full-operator and BDDC-specific branches stay inactive.
\end{itemize}

This is the cleanest summary of what Newton is actually doing on the default
parallel PMG route.

\section{Plain Newton On The Mainline Path}

One Newton iteration for the current benchmark is:
\begin{enumerate}
\item ask the constitutive builder for the current free-output residual and
tangent information,
\item form the regularized operator
\[
K_r = r K_{\mathrm{elast}} + (1-r) K_{\mathrm{tangent}},
\]
on the owned-row PETSc pattern,
\item extract or reuse the reduced free-space operator,
\item solve the correction with the PMG-shell-preconditioned Krylov solver,
\item run the line search,
\item update the state and regularization level.
\end{enumerate}

The nonlinear method is still Newton. The main change is the operator format.

\section{Nested SSR Solve On The Mainline Path}

Because the current benchmark is an indirect SSR solve, the active Newton
routine is \code{newton\_ind\_ssr(...)}. It still follows the classical MATLAB
structure:
\begin{itemize}
\item build the \(V\) correction associated with the current residual,
\item build the \(W\) sensitivity direction,
\item recover \(d_\lambda\) from the work constraint,
\item combine them into the displacement correction,
\item damp the coupled update through \code{damping\_alg5(...)}.
\end{itemize}

What differs from MATLAB is that both \(V\) and \(W\) are solved through the
same reduced PMG-shell linear-system interface.

\section{Stopping And Damping In The Current Benchmark}

The committed benchmark config uses:
\begin{itemize}
\item \code{stopping\_criterion = absolute\_delta\_lambda},
\item \code{stopping\_tol = 1e-4},
\item \code{it\_damp\_max = 10},
\item \code{r\_min = 1e-4}.
\end{itemize}

So the mainline nonlinear solve is tuned around SSR progress in
\(\lambda\), not around a pure residual-only stopping rule.

\section{What Newton Receives From The Constitutive Layer}

On the mainline path, Newton primarily sees:
\begin{itemize}
\item free-output residual builders,
\item owned-row tangent refreshes hidden behind the constitutive builder,
\item cached PETSc regularized matrices on the owned pattern.
\end{itemize}

This is why the tangent and Newton chapters have to be read together. The
current Newton branch is inseparable from the owned-row assembly path beneath
it.

\section{Deflation Basis On The Mainline Path}

The solver-side deflation basis still grows after useful accepted directions,
just as in MATLAB. On the current mainline path, that means:
\begin{itemize}
\item the basis is owned by the PETSc solver object,
\item continuation and Newton can reuse it across accepted steps,
\item the default benchmark does \textbf{not} enable temporary first-iteration
warm-start modes.
\end{itemize}

\section{Where To Edit}

\begin{itemize}
\item \petscmod{nonlinear/newton.py} for iteration flow and operator selection,
\item \petscmod{nonlinear/damping.py} for line-search behavior,
\item \petscmod{constitutive/problem.py} for the residual and regularized-matrix
interfaces used by Newton,
\item \petscmod{linear/solver.py} for backend-specific expectations.
\end{itemize}

\section{Where The Inactive Newton Branches Went}

Full-system operator mode, local full-operator branches, separate explicit
preconditioning matrices, and optional warm-start features are moved to
Appendix~\ref{app:newton}.
